\section{Introduction}
Test-time entropy minimization changes a deployed classifier without reading target labels. Its objective can decrease while classification risk first improves and then rises. This deterioration is established prior work: logit-adjusted test-time adaptation analyzes drift under partial class imbalance and proposes a correction~\cite{weerasinghe2026}. Our question is quantitative: when does the actual stochastic recursion first become worse than the frozen classifier on the same target distribution?

Answering this question requires more than identifying an unstable equilibrium. Translation changes the initialization relative to the entropy peak, while prior shift moves the peak and the supervised optimum differently. Their effects can cancel, reversing exit direction and delaying risk crossing. The boundary of the risk-safe region matters as much as local instability. Moreover, increasing the batch size suppresses SGD noise without suppressing calibration error, so a predictor that appears close in parameter space can become less reliable.

We analyze a fixed-feature binary classifier whose scalar threshold adapts by its actual entropy gradient. This restriction makes the objective, risk, and noise law explicit. Our theoretical contribution is the connection between risk-defined boundaries and the nonlinear, state-dependent mini-batch recursion, retaining integer update times and distinguishing risk-tolerance regimes. A general scalar extension separates field calibration from risk-boundary estimation. Its critical calibration limit predicts persistent random forecast error when peak uncertainty is comparable to the SGD seed scale.

The empirical contribution is a sequence of tests with different information levels. Gaussian experiments test mixed-shift cancellation and tolerance scaling. Frozen visual features test whether correct empirical geometry explains exits and whether independently estimated geometry predicts them. The latter initially fails. Target-geometry replacements localize the failure, and a larger independent calibration study recovers part of the gap. A remaining long-tail failure is investigated with modified-noise recursions, which support accumulated passage noise as a finite-parameter limitation of an early-seed approximation.

\paragraph{Claims and scope.}
We claim model-specific risk-exit laws, empirical evidence for the role of nonlinear geometry, and evidence separating calibration and continuing-noise errors. The strongest visual prediction uses labeled oracle risk endpoints. We therefore make no claim of a completed label-free stopping method, general BN adaptation, or population-risk certification. Proofs of the declared model statements are supplied in the supplement; their independent mathematical verification remains part of the research work.

\section{Relation to Prior Work}
\paragraph{Entropy drift and adaptation.}
Entropy minimization for test-time adaptation is exemplified by Tent~\cite{wang2021}. Weerasinghe et al.~\cite{weerasinghe2026} already establish imbalance-related degradation, normalization drift, and logit adjustment. Those phenomena are not contributions here. Their All-Class Accuracy and peak-to-average metric differ from our first crossing of the frozen classifier's risk on a fixed target distribution. A high probability of escaping toward the majority class is not a safety guarantee under our comparator.

\paragraph{Classical random exit.}
Logarithmic escape and random-seed corrections near unstable equilibria are classical~\cite{kifer1981,bakhtin2008}. Likewise, convolution closure of log-concavity is standard~\cite{saumard2014}. Our statements specialize these mechanisms to actual entropy gradients, classification-risk endpoints, discrete time, and independent geometry estimation. Whether these refinements constitute sufficient novelty for a vision contribution requires a direct comparison beyond the numerical evidence. We do not claim a new general instability mechanism.
